% --- Archon physics formalization source begin ---
% archon:physics
% archon:covers PhyXMiniProblems/problem_phyx_mini_0843.lean
% archon:source-report reports/phyx_mini/problem_phyx_mini_0843.source.json

\chapter{Physics problem phyx\_mini\_0843}
\label{ch:PhyXMiniProblems_problem_phyx_mini_0843}

\paragraph{Problem source.}
\#\# Question

What is the electric flux through the surface?

\#\# Answer choices

- A: A: 1.2N m\textasciicircum{}2/C

- B: B: 10N m\textasciicircum{}2/C

- C: C: 2N m\textasciicircum{}2/C

- D: D: 1N m\textasciicircum{}2/C

\#\# Figure caption (auxiliary; use the image as primary evidence)

The image depicts a rectangular plane with dimensions labeled as "10 cm × 10 cm." Several parallel arrows in magenta are shown passing through the plane at an angle. These arrows represent an electric field, labeled as \textbackslash{}( \textbackslash{}vec\{E\} \textbackslash{}), with a magnitude of "200 N/C."

A single black arrow labeled as \textbackslash{}( \textbackslash{}hat\{n\} \textbackslash{}) is perpendicular to the plane, indicating the normal vector of the surface. The angle between the plane and the direction of the electric field is marked as "30°."

\#\# Dataset metadata

Electrostatics; Physical Model Grounding Reasoning, Multi-Formula Reasoning

\paragraph{Recorded answer/context.}
D: D: 1N m\textasciicircum{}2/C

\paragraph{Figure/image path.}
/root/proposal\_for\_physic/science-mango/phyx\_mini\_run/phyx\_data/test\_image/843.png

\paragraph{Formalization target.}
create a compiling Lean file with sorry bodies at `PhyXMiniProblems/problem\_phyx\_mini\_0843.lean`. The Lean declarations must preserve the physical quantities, dimensions or dimensional roles, figure labels, governing-law hypotheses, and final relation expressed by this problem.
Use Mathlib/Physlib names found through LeanExplore where available. If a physics API is missing, introduce faithful local abstractions rather than scalar placeholder aliases.

\begin{theorem}[Physics formalization target]
\label{thm:physics:phyx_mini_0843:target}
\lean{PhyXMiniProblems.ProblemPhyXMini0843.problem_phyx_mini_0843}
\uses{def:physics:phyx-mini-0843:phyxminiproblems-problemphyxmini0843-electricfluxinnewtonmetressquaredpercoulomb, def:physics:phyx-mini-0843:phyxminiproblems-problemphyxmini0843-rectangularsurfacefluxsetup, def:physics:phyx-mini-0843:phyxminiproblems-problemphyxmini0843-matchessuppliedfigure, def:physics:phyx-mini-0843:phyxminiproblems-problemphyxmini0843-satisfiesfieldplaneanglegeometry, def:physics:phyx-mini-0843:phyxminiproblems-problemphyxmini0843-satisfiesuniformrectangularelectricfluxlaw, def:physics:phyx-mini-0843:phyxminiproblems-problemphyxmini0843-displayedelectricfluxinnewtonmetressquaredpercoulomb}
The electric flux through the \(10\,\mathrm{cm}\times10\,\mathrm{cm}\)
surface is \(1\,\mathrm{N\,m^2/C}\), answer D.
\end{theorem}
\begin{proof}
The field makes \(30^\circ\) with the plane, so its normal component is
\(E\sin30^\circ=200/2=100\,\mathrm{N/C}\).  The area is
\((0.10\,\mathrm m)^2=0.01\,\mathrm{m^2}\).  The uniform-field flux law gives
\[
 \Phi_E=0.01(100)=1\,\mathrm{N\,m^2/C},
\]
which is the value displayed by D.
\end{proof}

% --- Archon declaration topology begin ---
\section{Declaration topology}
\begin{definition}[electric Field Dimension]
  \label{def:physics:phyx-mini-0843:phyxminiproblems-problemphyxmini0843-electricfielddimension}
  \lean{PhyXMiniProblems.ProblemPhyXMini0843.electricFieldDimension}
  The MLTQ dimension of force per charge, hence of an electric field
\end{definition}
\begin{proof}
  This is the declared model definition of electric Field Dimension.
\end{proof}
\begin{definition}[electric Flux Dimension]
  \label{def:physics:phyx-mini-0843:phyxminiproblems-problemphyxmini0843-electricfluxdimension}
  \lean{PhyXMiniProblems.ProblemPhyXMini0843.electricFluxDimension}
  \uses{def:physics:phyx-mini-0843:phyxminiproblems-problemphyxmini0843-electricfielddimension}
  Electric flux has dimension electric field times area
\end{definition}
\begin{proof}
  This is the declared model definition of electric Flux Dimension.
\end{proof}
\begin{definition}[Length Quantity]
  \label{def:physics:phyx-mini-0843:phyxminiproblems-problemphyxmini0843-lengthquantity}
  \lean{PhyXMiniProblems.ProblemPhyXMini0843.LengthQuantity}
  A nonnegative, unit-independent physical length
\end{definition}
\begin{proof}
  This is the declared model definition of Length Quantity.
\end{proof}
\begin{definition}[Electric Field Vector Quantity]
  \label{def:physics:phyx-mini-0843:phyxminiproblems-problemphyxmini0843-electricfieldvectorquantity}
  \lean{PhyXMiniProblems.ProblemPhyXMini0843.ElectricFieldVectorQuantity}
  \uses{def:physics:phyx-mini-0843:phyxminiproblems-problemphyxmini0843-electricfielddimension}
  A unit-independent three-dimensional electric-field vector
\end{definition}
\begin{proof}
  This is the declared model definition of Electric Field Vector Quantity.
\end{proof}
\begin{definition}[Electric Flux Quantity]
  \label{def:physics:phyx-mini-0843:phyxminiproblems-problemphyxmini0843-electricfluxquantity}
  \lean{PhyXMiniProblems.ProblemPhyXMini0843.ElectricFluxQuantity}
  \uses{def:physics:phyx-mini-0843:phyxminiproblems-problemphyxmini0843-electricfluxdimension}
  Signed electric flux through an oriented surface
\end{definition}
\begin{proof}
  This is the declared model definition of Electric Flux Quantity.
\end{proof}
\begin{definition}[length Readout]
  \label{def:physics:phyx-mini-0843:phyxminiproblems-problemphyxmini0843-lengthreadout}
  \lean{PhyXMiniProblems.ProblemPhyXMini0843.lengthReadout}
  \uses{def:physics:phyx-mini-0843:phyxminiproblems-problemphyxmini0843-lengthquantity}
  Scalar readout of a physical length in a coherent choice of units
\end{definition}
\begin{proof}
  This is the declared model definition of length Readout.
\end{proof}
\begin{definition}[electric Field Vector Readout]
  \label{def:physics:phyx-mini-0843:phyxminiproblems-problemphyxmini0843-electricfieldvectorreadout}
  \lean{PhyXMiniProblems.ProblemPhyXMini0843.electricFieldVectorReadout}
  \uses{def:physics:phyx-mini-0843:phyxminiproblems-problemphyxmini0843-electricfieldvectorquantity}
  Vector readout of an electric field in a coherent choice of units
\end{definition}
\begin{proof}
  This is the declared model definition of electric Field Vector Readout.
\end{proof}
\begin{definition}[electric Flux Readout]
  \label{def:physics:phyx-mini-0843:phyxminiproblems-problemphyxmini0843-electricfluxreadout}
  \lean{PhyXMiniProblems.ProblemPhyXMini0843.electricFluxReadout}
  \uses{def:physics:phyx-mini-0843:phyxminiproblems-problemphyxmini0843-electricfluxquantity}
  Scalar readout of electric flux in a coherent choice of units
\end{definition}
\begin{proof}
  This is the declared model definition of electric Flux Readout.
\end{proof}
\begin{definition}[length In Centimeters]
  \label{def:physics:phyx-mini-0843:phyxminiproblems-problemphyxmini0843-lengthincentimeters}
  \lean{PhyXMiniProblems.ProblemPhyXMini0843.lengthInCentimeters}
  \uses{def:physics:phyx-mini-0843:phyxminiproblems-problemphyxmini0843-lengthquantity, def:physics:phyx-mini-0843:phyxminiproblems-problemphyxmini0843-lengthreadout}
  Length readout in centimetres, matching the surface labels in the figure
\end{definition}
\begin{proof}
  This is the declared model definition of length In Centimeters.
\end{proof}
\begin{definition}[electric Field Vector In Newtons Per Coulomb]
  \label{def:physics:phyx-mini-0843:phyxminiproblems-problemphyxmini0843-electricfieldvectorinnewtonspercoulomb}
  \lean{PhyXMiniProblems.ProblemPhyXMini0843.electricFieldVectorInNewtonsPerCoulomb}
  \uses{def:physics:phyx-mini-0843:phyxminiproblems-problemphyxmini0843-electricfieldvectorquantity, def:physics:phyx-mini-0843:phyxminiproblems-problemphyxmini0843-electricfieldvectorreadout}
  Electric-field vector readout in SI units, i.e. newtons per coulomb
\end{definition}
\begin{proof}
  This is the declared model definition of electric Field Vector In Newtons Per Coulomb.
\end{proof}
\begin{definition}[electric Flux In Newton Metres Squared Per Coulomb]
  \label{def:physics:phyx-mini-0843:phyxminiproblems-problemphyxmini0843-electricfluxinnewtonmetressquaredpercoulomb}
  \lean{PhyXMiniProblems.ProblemPhyXMini0843.electricFluxInNewtonMetresSquaredPerCoulomb}
  \uses{def:physics:phyx-mini-0843:phyxminiproblems-problemphyxmini0843-electricfluxquantity, def:physics:phyx-mini-0843:phyxminiproblems-problemphyxmini0843-electricfluxreadout}
  Electric-flux readout in SI units, i.e. newton-metres squared per coulomb
\end{definition}
\begin{proof}
  This is the declared model definition of electric Flux In Newton Metres Squared Per Coulomb.
\end{proof}
\begin{definition}[degrees To Radians]
  \label{def:physics:phyx-mini-0843:phyxminiproblems-problemphyxmini0843-degreestoradians}
  \lean{PhyXMiniProblems.ProblemPhyXMini0843.degreesToRadians}
  Convert a numerical angle in degrees to its radian readout
\end{definition}
\begin{proof}
  This is the declared model definition of degrees To Radians.
\end{proof}
\begin{definition}[Oriented Rectangular Surface]
  \label{def:physics:phyx-mini-0843:phyxminiproblems-problemphyxmini0843-orientedrectangularsurface}
  \lean{PhyXMiniProblems.ProblemPhyXMini0843.OrientedRectangularSurface}
  \uses{def:physics:phyx-mini-0843:phyxminiproblems-problemphyxmini0843-lengthquantity}
  An oriented rectangular plane. The two physical side lengths determine its area, while unitNormal is the direction denoted by n-hat in the figure
\end{definition}
\begin{proof}
  This is the declared model definition of Oriented Rectangular Surface.
\end{proof}
\begin{definition}[Rectangular Surface Flux Setup]
  \label{def:physics:phyx-mini-0843:phyxminiproblems-problemphyxmini0843-rectangularsurfacefluxsetup}
  \lean{PhyXMiniProblems.ProblemPhyXMini0843.RectangularSurfaceFluxSetup}
  \uses{def:physics:phyx-mini-0843:phyxminiproblems-problemphyxmini0843-electricfieldvectorquantity, def:physics:phyx-mini-0843:phyxminiproblems-problemphyxmini0843-electricfluxquantity, def:physics:phyx-mini-0843:phyxminiproblems-problemphyxmini0843-orientedrectangularsurface}
  The physical quantities appearing in the diagram. electricFlux is an independent signed physical quantity; its relation to the other fields is imposed only by the governing flux law below
\end{definition}
\begin{proof}
  This is the declared model definition of Rectangular Surface Flux Setup.
\end{proof}
\begin{definition}[Matches Supplied Figure]
  \label{def:physics:phyx-mini-0843:phyxminiproblems-problemphyxmini0843-matchessuppliedfigure}
  \lean{PhyXMiniProblems.ProblemPhyXMini0843.MatchesSuppliedFigure}
  \uses{def:physics:phyx-mini-0843:phyxminiproblems-problemphyxmini0843-lengthincentimeters, def:physics:phyx-mini-0843:phyxminiproblems-problemphyxmini0843-electricfieldvectorinnewtonspercoulomb, def:physics:phyx-mini-0843:phyxminiproblems-problemphyxmini0843-degreestoradians, def:physics:phyx-mini-0843:phyxminiproblems-problemphyxmini0843-rectangularsurfacefluxsetup}
  Direct transcription of the primary image: both side labels are 10 cm, the field-magnitude label is 200 N/C, the field-plane angle is 30 degrees, and the hatted normal is a unit vector. No electric-flux value or answer choice occurs in this structure
\end{definition}
\begin{proof}
  This is the declared model definition of Matches Supplied Figure.
\end{proof}
\begin{definition}[Satisfies Field Plane Angle Geometry]
  \label{def:physics:phyx-mini-0843:phyxminiproblems-problemphyxmini0843-satisfiesfieldplaneanglegeometry}
  \lean{PhyXMiniProblems.ProblemPhyXMini0843.SatisfiesFieldPlaneAngleGeometry}
  \uses{def:physics:phyx-mini-0843:phyxminiproblems-problemphyxmini0843-electricfieldvectorreadout, def:physics:phyx-mini-0843:phyxminiproblems-problemphyxmini0843-rectangularsurfacefluxsetup}
  For the acute, positively oriented angle shown between the field and the plane, the field component along the surface normal is |E| sin(theta). This is stated in every coherent unit choice and contains no numerical answer
\end{definition}
\begin{proof}
  This is the declared model definition of Satisfies Field Plane Angle Geometry.
\end{proof}
\begin{definition}[Satisfies Uniform Rectangular Electric Flux Law]
  \label{def:physics:phyx-mini-0843:phyxminiproblems-problemphyxmini0843-satisfiesuniformrectangularelectricfluxlaw}
  \lean{PhyXMiniProblems.ProblemPhyXMini0843.SatisfiesUniformRectangularElectricFluxLaw}
  \uses{def:physics:phyx-mini-0843:phyxminiproblems-problemphyxmini0843-lengthreadout, def:physics:phyx-mini-0843:phyxminiproblems-problemphyxmini0843-electricfieldvectorreadout, def:physics:phyx-mini-0843:phyxminiproblems-problemphyxmini0843-electricfluxreadout, def:physics:phyx-mini-0843:phyxminiproblems-problemphyxmini0843-rectangularsurfacefluxsetup}
  The uniform-field electric-flux law for an oriented rectangular plane: Phi\_E = area * (E dot n-hat). It is quantified over coherent unit choices, so the equality retains its dimensional meaning and mentions no answer value
\end{definition}
\begin{proof}
  This is the declared model definition of Satisfies Uniform Rectangular Electric Flux Law.
\end{proof}
\begin{definition}[Answer Choice]
  \label{def:physics:phyx-mini-0843:phyxminiproblems-problemphyxmini0843-answerchoice}
  \lean{PhyXMiniProblems.ProblemPhyXMini0843.AnswerChoice}
  Labels of the four answer choices printed with the problem
\end{definition}
\begin{proof}
  This is the declared model definition of Answer Choice.
\end{proof}
\begin{definition}[displayed Electric Flux In Newton Metres Squared Per Coulomb]
  \label{def:physics:phyx-mini-0843:phyxminiproblems-problemphyxmini0843-displayedelectricfluxinnewtonmetressquaredpercoulomb}
  \lean{PhyXMiniProblems.ProblemPhyXMini0843.displayedElectricFluxInNewtonMetresSquaredPerCoulomb}
  \uses{def:physics:phyx-mini-0843:phyxminiproblems-problemphyxmini0843-answerchoice}
  Numerical SI electric-flux readout printed beside an answer choice, in newton-metres squared per coulomb
\end{definition}
\begin{proof}
  This is the declared model definition of displayed Electric Flux In Newton Metres Squared Per Coulomb.
\end{proof}
% --- Archon declaration topology end ---
% --- Archon physics formalization source end ---
